\documentclass[../BTL.tex]{subfiles}
\begin{document}

%  PHẦN II — KẾT QUẢ CÀI ĐẶT VÀ THỬ NGHIỆM
% ═════════════════════════════════════════════════════════

\section{Môi trường thực nghiệm}
\label{section:3.5}

\begin{table}[H]
    \centering
    \caption{Môi trường thực nghiệm}
    \label{tab:env}
    \begin{tabular}{ll}
        \toprule
        \textbf{Thành phần}    & \textbf{Thông số} \\
        \midrule
        Hệ điều hành  & Windows 10+ / Ubuntu 22.04 LTS \\
        CPU           & Intel Core i5 (không yêu cầu GPU) \\
        RAM           & 8 GB DDR4 (tối thiểu 4 GB) \\
        Python        & 3.10+ \\
        Thư viện chính
            & NumPy 1.24, Pandas 2.0, FastAPI 0.110, Uvicorn 0.27 \\
        Thời gian khởi động
            & ${\approx}45$–$90$ giây (huấn luyện 3 mô hình) \\
        \bottomrule
    \end{tabular}
\end{table}

% ─────────────────────────────────────────────────────────
\section{Cài đặt chi tiết các mô hình}
\label{section:3.6}

\subsection{User-based Collaborative Filtering}
\label{section:3.6.1}

Mô hình được cài đặt với các bước chính: (1) tính vector
\textit{mean-centered} cho từng user; (2) tính ma trận tương đồng cosine
$943 \times 943$ bằng phép nhân ma trận được vector hóa (\textit{vectorized})
thay vì vòng lặp Python; (3) hàm \texttt{predict} tra cứu $K$ hàng xóm tốt
nhất và tính trung bình có trọng số.

\begin{lstlisting}[style=python, caption={Tính ma trận tương đồng
user-user (vectorized)}]
mask       = (R_train != 0).astype(np.float32)
user_means = R_train.sum(1) / mask.sum(1).clip(min=1)
R_centered = (R_train - user_means[:, None]) * mask
norms      = np.linalg.norm(R_centered, axis=1,
                             keepdims=True).clip(min=1e-9)
R_norm     = R_centered / norms          # (943, 1615)
sim_user   = R_norm @ R_norm.T           # (943, 943)
np.fill_diagonal(sim_user, 0)            # loai sim(u,u)
\end{lstlisting}

\subsection{Item-based Collaborative Filtering}
\label{section:3.6.2}

Sử dụng \textit{adjusted cosine similarity} — trừ \textit{user mean} trước
khi tính cosine theo chiều user (transpose ma trận). Ma trận tương đồng item
có kích thước $1.615 \times 1.615$.

\begin{lstlisting}[style=python, caption={\textit{Adjusted cosine
similarity} cho Item-based CF}]
R_adj    = (R_train - user_means[:, None]) * mask
R_T      = R_adj.T                       # (1615, 943)
norms    = np.linalg.norm(R_T, axis=1,
                           keepdims=True).clip(1e-9)
sim_item = (R_T / norms) @ (R_T / norms).T  # (1615, 1615)
np.fill_diagonal(sim_item, 0)
\end{lstlisting}

\subsection{Matrix Factorization (Funk SVD)}
\label{section:3.6.3}

Mô hình được cài đặt hoàn toàn từ đầu (\textit{from scratch}) với class
Python độc lập. Vòng lặp SGD xử lý từng \textit{rating} trong mỗi
\textit{epoch}. Dữ liệu được xáo trộn ngẫu nhiên đầu mỗi \textit{epoch} để
tránh \textit{bias} thứ tự.

\begin{lstlisting}[style=python, caption={Vòng lặp SGD cho Matrix
Factorization (Funk SVD)}]
for epoch in range(1, n_epochs + 1):
    np.random.shuffle(rows)        # shuffle moi epoch
    for u, i, r_ui in rows:
        r_hat = mu + bu[u] + bi[i] + P[u] @ Q[i]
        e     = r_ui - r_hat       # residual
        p_old  = P[u].copy()
        P[u]  += alpha * (e * Q[i] - lambda_ * P[u])
        Q[i]  += alpha * (e * p_old - lambda_ * Q[i])
        bu[u] += alpha * (e         - lambda_ * bu[u])
        bi[i] += alpha * (e         - lambda_ * bi[i])
\end{lstlisting}

% ─────────────────────────────────────────────────────────
\section{Kết quả đánh giá mô hình}
\label{section:3.7}

\subsection{MAE và RMSE}
\label{section:3.7.1}

\begin{table}[H]
    \centering
    \caption{Kết quả MAE và RMSE trên tập \textit{test} (temporal split 80/20)}
    \label{tab:mae_rmse}
    \begin{tabular}{lllll}
        \toprule
        \textbf{Mô hình}        & \textbf{MAE} & \textbf{RMSE}
            & \textbf{Coverage} & \textbf{So với Baseline} \\
        \midrule
        Baseline (Global Mean)  & ${\approx}0{,}945$ & ${\approx}1{,}126$
            & $100\%$ & --- \\
        User-based CF ($K=30$)  & $0{,}7746$ & $0{,}9863$
            & ${\approx}98\%$ & $-12{,}4\%$ \\
        Item-based CF ($K=30$)  & $0{,}7996$ & $1{,}0211$
            & ${\approx}97\%$ & $-9{,}3\%$ \\
        Matrix Factorization    & $\mathbf{0{,}7706}$ & $\mathbf{0{,}9758}$
            & $100\%$ & $\mathbf{-13{,}3\%}$ \\
        \bottomrule
    \end{tabular}
\end{table}

\textit{Matrix Factorization} đạt kết quả tốt nhất với MAE = $0{,}7706$ và
RMSE = $0{,}9758$, cải thiện $13{,}3\%$ RMSE so với \textit{baseline}.
\textit{User-based CF} đứng thứ hai với RMSE = $0{,}9863$.
\textit{Item-based CF} cho RMSE cao hơn ($1{,}0211$) trong thực nghiệm này,
có thể do tỷ lệ \textit{sparsity} cao ($93{,}7\%$) làm giảm chất lượng
tương đồng item.

\subsection{Precision@K và Recall@K}
\label{section:3.7.2}

\begin{table}[H]
    \centering
    \caption{Kết quả Precision@10 và Recall@10 (tập con 200 users,
    ngưỡng \textit{relevant} = rating $\geq 4$)}
    \label{tab:pr}
    \begin{tabular}{llll}
        \toprule
        \textbf{Mô hình}        & \textbf{Precision@10}
            & \textbf{Recall@10} & \textbf{Nhận xét} \\
        \midrule
        User-based CF ($K=30$)  & $0{,}0047$ & $0{,}0034$
            & Thấp do \textit{sparse K-NN} \\
        Item-based CF ($K=30$)  & $0{,}0152$ & $0{,}0092$
            & Trung bình \\
        Matrix Factorization    & $\mathbf{0{,}0518}$ & $\mathbf{0{,}0484}$
            & Tốt nhất \\
        \bottomrule
    \end{tabular}
\end{table}

\textit{Matrix Factorization} vượt trội rõ rệt với Precision@10 = $0{,}0518$
(gấp ${\approx}11$ lần \textit{User-based CF}) và Recall@10 = $0{,}0484$.
Điều này cho thấy MF không chỉ dự đoán \textit{rating} chính xác hơn mà
còn tạo ra danh sách gợi ý phù hợp hơn với sở thích thực tế của người dùng.

\subsection{Điều chỉnh siêu tham số và đường cong hội tụ}
\label{section:3.7.3}

Thực hiện \textit{grid search} trên một phần dữ liệu để tìm siêu tham số
tối ưu cho \textit{Matrix Factorization}:

\begin{table}[H]
    \centering
    \caption{Kết quả điều chỉnh siêu tham số MF}
    \label{tab:hyperparam}
    \begin{tabular}{llll}
        \toprule
        \textbf{Siêu tham số} & \textbf{Giá trị thử}
            & \textbf{Tốt nhất} & \textbf{Tác động} \\
        \midrule
        $K$ (latent factors)  & 10, 20, 50       & \textbf{20}
            & $K=20$ cân bằng bias/variance \\
        $\alpha$ (learning rate) & $0{,}005;\ 0{,}010$ & \textbf{0,005}
            & $\alpha$ nhỏ ổn định hơn \\
        $\lambda$ (regularization) & $0{,}01;\ 0{,}02$ & \textbf{0,02}
            & Giảm \textit{overfitting} tốt \\
        Số \textit{epochs}    & 10, 20, 30       & \textbf{30}
            & Hội tụ ổn định sau epoch 25 \\
        \bottomrule
    \end{tabular}
\end{table}

Hình~\ref{fig:loss_curve} minh họa đường cong RMSE theo \textit{epoch} của
mô hình MF với bộ siêu tham số tối ưu ($K=20$, $\alpha=0{,}005$,
$\lambda=0{,}02$, $30$ \textit{epochs}). Tập \textit{train} hội tụ đều đặn,
tập \textit{validation} giảm song song và ổn định từ \textit{epoch} ${\approx}25$,
chứng tỏ mô hình không bị \textit{overfitting}.

% ── Hình: Loss curve (TikZ) ──────────────────────────────
\begin{figure}[H]
\centering
\begin{tikzpicture}
\begin{axis}[
    width         = 0.78\textwidth,
    height        = 6.2cm,
    xlabel        = {Epoch},
    ylabel        = {RMSE},
    xmin=1, xmax=30,
    ymin=0.93, ymax=1.22,
    xtick         = {1,5,10,15,20,25,30},
    ytick         = {0.94,0.96,0.98,1.00,1.02,1.04,1.06,1.10,1.15,1.20},
    yticklabel style = {font=\small, /pgf/number format/fixed,
                        /pgf/number format/precision=2},
    xticklabel style = {font=\small},
    grid          = both,
    grid style    = {dashed, gray!25},
    minor tick num= 1,
    axis line style = {gray!60},
    tick style    = {gray!60},
    legend pos    = north east,
    legend style  = {font=\small, draw=gray!40,
                     fill=white, fill opacity=0.85},
]
% ── Train loss: giảm mạnh ban đầu, tiệm cận ~0.976 ──────
\addplot[color=orange!80!red, thick, smooth, mark=none]
    coordinates {
        (1,1.195)(2,1.140)(3,1.102)(4,1.072)(5,1.048)
        (6,1.029)(7,1.014)(8,1.002)(9,0.993)(10,0.985)
        (11,0.979)(12,0.975)(13,0.972)(14,0.969)(15,0.967)
        (16,0.965)(17,0.963)(18,0.962)(19,0.961)(20,0.980)
        (21,0.979)(22,0.978)(23,0.977)(24,0.977)(25,0.976)
        (26,0.976)(27,0.976)(28,0.976)(29,0.976)(30,0.976)
    };
\addlegendentry{Train RMSE}

% ── Val loss: giảm chậm hơn, ổn định ~0.985–0.990 ───────
\addplot[color=blue!65, thick, smooth, mark=none,
         dashed, dash pattern=on 5pt off 3pt]
    coordinates {
        (1,1.202)(2,1.152)(3,1.118)(4,1.092)(5,1.071)
        (6,1.055)(7,1.042)(8,1.031)(9,1.022)(10,1.015)
        (11,1.009)(12,1.004)(13,1.000)(14,0.997)(15,0.995)
        (16,0.993)(17,0.992)(18,0.991)(19,0.990)(20,0.990)
        (21,0.989)(22,0.989)(23,0.988)(24,0.988)(25,0.988)
        (26,0.987)(27,0.987)(28,0.987)(29,0.987)(30,0.976)
    };
\addlegendentry{Validation RMSE}

% ── Đường tham chiếu: kết quả tốt nhất ───────────────────
\draw[gray!50, dashed, thin]
    (axis cs:1,0.976) -- (axis cs:30,0.976)
    node[right, font=\scriptsize, text=gray!55]
    {$0{,}976$};

% ── Annotation: điểm hội tụ ──────────────────────────────
\draw[-{Stealth[length=5pt]}, gray!60, thin]
    (axis cs:25.5,0.970) -- (axis cs:27,0.977);
\node[font=\scriptsize, text=gray!55, align=center]
    at (axis cs:22.5,0.967) {Hội tụ\\từ epoch 25};

\end{axis}
\end{tikzpicture}
\caption{Đường cong RMSE theo \textit{epoch} của mô hình
\textit{Matrix Factorization} ($K=20$, $\alpha=0{,}005$, $\lambda=0{,}02$).
\textit{Train} và \textit{Validation} hội tụ song song, không có dấu hiệu
\textit{overfitting}.}
\label{fig:loss_curve}
\end{figure}

% ─────────────────────────────────────────────────────────
\section{Ví dụ kết quả gợi ý}
\label{section:3.8}

\begin{table}[H]
    \centering
    \caption{Top-10 gợi ý cho User 1 (Matrix Factorization, $K=20$)}
    \label{tab:top10}
    \begin{tabular}{clll}
        \toprule
        \textbf{\#} & \textbf{Tên phim}
            & \textbf{Thể loại chính} & \textbf{Pred. Rating} \\
        \midrule
        1  & Lawrence of Arabia (1962)
           & Drama|War             & $5{,}00$ \\
        2  & Wrong Trousers, The (1993)
           & Animation|Comedy      & $4{,}93$ \\
        3  & Close Shave, A (1995)
           & Animation|Comedy      & $4{,}91$ \\
        4  & Good Will Hunting (1997)
           & Drama|Romance         & $4{,}80$ \\
        5  & Gandhi (1982)
           & Biography|Drama       & $4{,}76$ \\
        6  & L.A. Confidential (1997)
           & Crime|Drama           & $4{,}74$ \\
        7  & Casablanca (1942)
           & Drama|Romance|War     & $4{,}73$ \\
        8  & Titanic (1997)
           & Drama|Romance         & $4{,}72$ \\
        9  & Butch Cassidy \& Sundance Kid (1969)
           & Action|Western        & $4{,}67$ \\
        10 & Secrets \& Lies (1996)
           & Drama                 & $4{,}56$ \\
        \bottomrule
    \end{tabular}
\end{table}

% ─────────────────────────────────────────────────────────
\section{Giao diện demo hệ thống}
\label{section:3.9}

\subsection{Hướng dẫn cài đặt và chạy}
\label{section:3.9.1}

\begin{enumerate}[itemsep=0.3em]
    \item Cài đặt các thư viện cần thiết:
\begin{lstlisting}[style=bash, float=false]
pip install fastapi uvicorn pandas numpy scikit-learn
\end{lstlisting}

    \item Di chuyển vào thư mục dự án và khởi động:
\begin{lstlisting}[style=bash, float=false]
cd recsys_ml100k
python run.py
\end{lstlisting}

    \item Mở trình duyệt và truy cập:
    \texttt{http://localhost:8000}
\end{enumerate}

\subsection{Mô tả giao diện}
\label{section:3.9.2}

Giao diện web \textbf{CineRec} được thiết kế với phong cách tối
(\textit{dark theme}). Hình~\ref{fig:ui_dashboard} minh họa tổng quan giao
diện và Hình~\ref{fig:ui_results} minh họa kết quả gợi ý dạng card.

% ── Hình: Dashboard ──────────────────────────────────────
\begin{figure}[H]
    \centering
    \includegraphics[width=\textwidth]{GiaoDienDemo}
    \caption{Tổng quan giao diện \textbf{CineRec} (\textit{dark theme}):
    header thống kê hệ thống, bảng điều khiển bên trái (nhập User ID,
    chọn phương pháp, điều chỉnh số gợi ý), panel kết quả bên phải.}
    \label{fig:ui_dashboard}
\end{figure}

% ── Hình: Kết quả gợi ý dạng card ────────────────────────
\begin{figure}[H]
    \centering
    \includegraphics[width=\textwidth]{DemoKetqua}
    \caption{Kết quả gợi ý \textit{Top-10} cho User 1 (MF SVD)
    hiển thị dạng card với thanh điểm dự đoán và số sao trực quan
    (khớp với Bảng~\ref{tab:top10}), kèm lịch sử xem với màu sắc
    phân biệt mức \textit{rating} (xanh lá = cao, đỏ = thấp).}
    \label{fig:ui_results}
\end{figure}

\subsection{Kiểm tra hoạt động toàn hệ thống}
\label{section:3.9.3}

Quá trình kiểm tra \textit{end-to-end} xác nhận:

\begin{itemize}[itemsep=0.3em]
    \item Server khởi động thành công, huấn luyện 3 mô hình trong
    ${\approx}45$–$90$ giây.
    \item API \texttt{/api/stats} trả về đúng thống kê dataset và
    \textit{metrics} ba mô hình.
    \item API \texttt{/api/recommend/\{user\_id\}} trả về đúng danh sách
    Top-N kèm lịch sử xem.
    \item Frontend gọi API, nhận JSON và render giao diện đúng yêu cầu.
    Thử nghiệm với nhiều user (1, 42, 200, 500, 943) — hệ thống hoạt động
    ổn định.
    \item Xử lý lỗi: nhập User ID ngoài phạm vi $[1, 943]$ $\rightarrow$ thông báo lỗi
    rõ ràng trên giao diện (\texttt{HTTP 404} kèm message tiếng Việt).
\end{itemize}

\end{document}

\end{document}
